\documentclass[journal]{IEEEtran}
\usepackage{fontspec}
\setmainfont{TeX Gyre Termes}
\usepackage{amsmath,amssymb}
\usepackage{graphicx}
\usepackage{booktabs}
\usepackage{cite}

\begin{document}
\title{Zone-Based Reliability Assessment of\\ Reconfigurable Distribution Networks}
\author{Alice~Zhang and Bob~Li}
\maketitle

\begin{abstract}
Reliability assessment of distribution networks usually assumes a fixed topology, while
operators reconfigure feeders and transfer loads many times a day. This paper groups the
buses between switching devices into zones whose internal structure never changes, and
expresses the supply relations between zones as a sparse matrix. Reliability indices of every
load point then follow from matrix products, without a new search of the network after each
switching operation. Tests on the IEEE 33-bus system and on interconnected feeders show that the
indices agree with an enumeration method while the computation time drops by an order of
magnitude.
\end{abstract}

\begin{IEEEkeywords}
Distribution network, reliability assessment, network reconfiguration, switching zones.
\end{IEEEkeywords}

\section*{Nomenclature}

\subsubsection*{Sets and Indices}
\begin{IEEEdescription}[\IEEEusemathlabelsep\IEEEsetlabelwidth{$\mathcal{Z},\mathcal{K}$}]
\item[$\mathcal{N},\mathcal{L}$] Sets of buses and of branches.
\item[$\mathcal{Z},\mathcal{K}$] Sets of zones and of switching devices.
\item[$z,k,j$] Indices of zones, switching devices and load points.
\end{IEEEdescription}

\subsubsection*{Parameters}
\begin{IEEEdescription}[\IEEEusemathlabelsep\IEEEsetlabelwidth{$\mathcal{Z},\mathcal{K}$}]
\item[$\lambda_z$] Failure rate of the branches in zone $z$, in failures per year.
\item[$r_z$] Mean repair time of a failure in zone $z$, in hours.
\item[$N_j$] Number of customers at load point $j$.
\end{IEEEdescription}

\subsubsection*{Variables}
\begin{IEEEdescription}[\IEEEusemathlabelsep\IEEEsetlabelwidth{$\mathcal{Z},\mathcal{K}$}]
\item[$\mathbf{B}$] Supply matrix; $B_{kz}=1$ when device $k$ lies on the path that feeds
  zone $z$.
\item[$U_j$] Annual unavailability of load point $j$, in hours per year.
\end{IEEEdescription}

\section{Introduction}
Distribution networks are planned as meshed structures but operated radially. Normally open tie
switches allow the operator to move loads between feeders, and a fault is cleared by opening the
nearest protective device and closing a tie switch so that healthy sections are restored. Each
of these operations changes which devices a customer depends on, and therefore its reliability.

Analytical methods evaluate the failure modes of each component and their effect on every load
point~\cite{billinton1996}. They are accurate, but they rebuild the upstream and downstream
relations of all devices whenever the topology changes. Simulation methods sample the operating
history instead~\cite{allan1991}, which suits planning studies better than the short horizon of
operation.

This paper observes that the buses between two switching devices always share the same supply
path. We call such a group a zone and show that the reliability indices of all load points can
be written in terms of a matrix of zone-to-device relations. The contributions are as follows.
First, the network is reduced to its zones once, before operation starts. Second, the supply
matrix is updated by local operations when a switch changes state. Third, transfers over tie
lines are included in the same matrix expressions.

\section{Zone Model}
\subsection{Zones and Supply Paths}
A zone is a maximal set of buses connected by branches that contain no switching device.
Because no branch inside a zone can be opened, all buses of a zone are supplied through the same
sequence of devices. Let $\mathcal{Z}$ be the set of zones and $\mathcal{K}$ the set of
switching devices. The supply matrix $\mathbf{B}$ has one row per device and one column per
zone, with
\begin{equation}
B_{kz} =
\begin{cases}
1, & \text{if device $k$ lies on the supply path of zone $z$,}\\
0, & \text{otherwise.}
\end{cases}
\label{eq:supply}
\end{equation}
For a radial configuration each column of $\mathbf{B}$ describes a path from a substation to the
zone, so the number of nonzero entries grows with the depth of the feeder rather than with its
number of buses.

\subsubsection{Fault Isolation}
When a branch in zone $z$ fails, the nearest upstream protective device opens and every zone it
feeds is interrupted. The zones downstream of the isolating switch are then restored through a
tie switch when one is available, so their outage lasts only the switching time.

\subsubsection{Restoration Order}
Zones upstream of the fault are restored first by reclosing the protective device after the
faulted zone has been isolated. The order does not change the indices of the faulted zone, whose
customers wait for the repair in every case.

\subsection{Reliability Indices}
With the supply matrix the annual unavailability of load point $j$ in zone $z(j)$ is
\begin{align}
U_j &= \sum_{z \in \mathcal{Z}} \lambda_z\, t_{z,z(j)}, \label{eq:unavail}\\
t_{z,z'} &=
\begin{cases}
r_z, & z' = z,\\
s, & z' \text{ restored through a tie switch},
\end{cases}\nonumber
\end{align}
where $s$ is the switching time. The system average interruption duration index follows from
\eqref{eq:unavail} as
\begin{equation}
\mathrm{SAIDI} = \frac{\sum_j N_j U_j}{\sum_j N_j}.
\label{eq:saidi}
\end{equation}
Only the columns of $\mathbf{B}$ that belong to zones affected by a switching operation change,
so the indices are updated in time proportional to the size of those zones.

\begin{table}[!t]
\caption{Test Systems}
\label{tab:systems}
\centering
\begin{tabular}{lrrr}
\toprule
System & Buses & Zones & Tie switches\\
\midrule
IEEE 33-bus & 33 & 8 & 5\\
Two feeders & 66 & 15 & 7\\
Three feeders & 135 & 24 & 11\\
\bottomrule
\end{tabular}
\end{table}

\section{Case Studies}
The method was applied to the systems listed in Table~\ref{tab:systems}. Failure rates and
repair times follow the data of~\cite{allan1991}, and every load point has the same number of
customers. The switching time of tie switches is one hour.

\begin{figure}[!t]
\centering
\rule{0.8\columnwidth}{0.4\columnwidth}
\caption{Zones of the IEEE 33-bus system. Each shaded block is a zone; tie switches are drawn
dashed.}
\label{fig:zones}
\end{figure}

\subsubsection{IEEE 33-Bus System}
Fig.~\ref{fig:zones} shows the eight zones of the IEEE 33-bus system. The supply matrix has
eight columns and twelve rows, against thirty-three buses in the original network. The SAIDI
computed from \eqref{eq:saidi} equals the value of the enumeration method to four digits.

\subsubsection{Interconnected Feeders}
For the three-feeder system the method evaluates all configurations reachable by one switching
operation in a fraction of the time the enumeration method needs for a single one. The gain grows
with the number of buses per zone, since the matrix size depends on the zones only.

Larger feeders keep the same structure: long laterals without switches become single zones, and
the work per operation stays small. This makes the indices available after every switching
operation, which is what operational reliability assessment requires.

\section{Conclusion}
Grouping buses into zones delimited by switching devices turns reliability assessment of a
reconfigurable distribution network into operations on a small supply matrix. The indices agree
with an enumeration method and are obtained much faster, so they can follow the network through
every reconfiguration. Future work will extend the zones to networks with distributed generation
that can supply islands during outages.

\begin{thebibliography}{1}
\bibitem{billinton1996}
R.~Billinton and R.~N. Allan, \emph{Reliability Evaluation of Power Systems}, 2nd~ed.
New York, NY, USA: Plenum, 1996.
\bibitem{allan1991}
R.~N. Allan, R.~Billinton, I.~Sjarief, L.~Goel, and K.~S. So, ``A reliability test system for
educational purposes---basic distribution system data and results,'' \emph{IEEE Trans. Power
Syst.}, vol.~6, no.~2, pp. 813--820, May 1991.
\end{thebibliography}
\end{document}
